% Paquetes usados
\usepackage{graphicx} % Imágenes
\usepackage{float} % imagenes donde las quieres [H]
\usepackage{amsmath} % Algunos símbolos matemáticos
\usepackage{amsfonts}
\usepackage{amssymb}
\usepackage{amsthm}
\usepackage{array} % Matrices, tablas
\usepackage{xcolor} % Colores de texto
\usepackage{enumitem} % Listas con letras
\setlist[itemize,1]{label=\raisebox{0.5ex}{\tiny$\bullet$}} % me gustan asi las listas
\usepackage[english]{babel} % inglés


% Distancias entre párrafos, quitar sangrías
\setlength{\parindent}{0pt}
\setlength{\parskip}{.8em}



% Título
\title{Lean's CIC Foundations}
\author{Pepa Montero}
\date{2026}

%clickable references
\usepackage[colorlinks=true, linkcolor=blue, citecolor=blue, urlcolor=blue]{hyperref}
\usepackage[all]{hypcap} % on clicked, goes to top of figure, not caption

%pastable references
% not been able to do it


% Theorem environments
\newtheorem{theorem}{Theorem}[section]
\newtheorem{definition}[theorem]{Definition}
\newtheorem{proposition}[theorem]{Proposition}
\newtheorem{lemma}[theorem]{Lemma}
\newtheorem{corollary}[theorem]{Corollary}
\newtheorem{theoremdef}[theorem]{Theorem/Definition}

\newtheoremstyle{notestyle} % name
  {6pt}                     % Space above
  {6pt}                     % Space below
  {\normalfont}                % Body font
  {}                        % Indent amount
  {\itshape}                % Head font (← italic instead of bold)
  {.}                       % Punctuation after head
  {0.5em}                   % Space after head
  {}                        % Head spec
\theoremstyle{notestyle}
\newtheorem*{note}{Note}
\AtBeginDocument{\renewcommand\proofname{Proof}}

% Box environment for reminders and stuff
\usepackage{tcolorbox}
\tcbuselibrary{theorems,breakable}
\newtcolorbox{highlight}[1][]{
  colback=white,              % background
  colframe=gray!50,           % light grey border
  boxrule=0.4pt,              % thinner border (default ~0.8pt)
  title=#1,
  fonttitle=\bfseries,
  breakable
}

% Computer Science pi lambda
\usepackage{pifont}
\makeatletter
\newcommand\Pimathsymbol[3][\mathord]{%
  #1{\@Pimathsymbol{#2}{#3}}}
\def\@Pimathsymbol#1#2{\mathchoice
  {\@Pim@thsymbol{#1}{#2}\tf@size}
  {\@Pim@thsymbol{#1}{#2}\tf@size}
  {\@Pim@thsymbol{#1}{#2}\sf@size}
  {\@Pim@thsymbol{#1}{#2}\ssf@size}}
\def\@Pim@thsymbol#1#2#3{%
  \mbox{\fontsize{#3}{#3}\Pisymbol{#1}{#2}}}
\makeatother
% the next two lines are needed to avoid LaTeX substituting upright from another font
\input{utxmia.fd}
\DeclareFontShape{U}{txmia}{m}{n}{<->ssub * txmia/m/it}{}
% you may also want
\DeclareFontShape{U}{txmia}{bx}{n}{<->ssub * txmia/bx/it}{}
% just in case
%\DeclareFontShape{U}{txmia}{l}{n}{<->ssub * txmia/l/it}{}
%\DeclareFontShape{U}{txmia}{b}{n}{<->ssub * txmia/b/it}{}
% plus info from Alan Munn at https://tex.stackexchange.com/questions/290165/how-do-i-get-a-nicer-lambda?noredirect=1#comment702120_290165
\newcommand{\pilambda}{\Pimathsymbol[\mathord]{txmia}{21}}


% Lean-like math commands
\newcommand{\real}{\mathbb R}
\newcommand{\R}{\mathbb R}
\newcommand{\C}{\mathbb C}
\newcommand{\Z}{\mathbb Z}

% Display style commands (like dfrac)
\newcommand{\dint}{\displaystyle\int}
\newcommand{\dsum}{\displaystyle\sum}

% Tikz pictures
\usepackage{tikz}
\usetikzlibrary{positioning, calc}
%\usetikzlibrary{arrows}
%\usetikzlibrary{arrows.meta}
%\usetikzlibrary{shapes.geometric}
\usepackage{tikz-cd}

% Código LEAN

\usepackage[T1]{fontenc}
\usepackage[utf8]{inputenc}

\usepackage{color}
\definecolor{keywordcolor}{rgb}{0.8, 0.1, 0.1}   % red
\definecolor{tacticcolor}{rgb}{0.0, 0.4, 0.9}    % blue
\definecolor{commentcolor}{rgb}{0.4, 0.4, 0.4}   % grey
\definecolor{symbolcolor}{rgb}{0.0, 0.4, 0.9}    % blue
\definecolor{sortcolor}{rgb}{0.1, 0.5, 0.1}      % green
\definecolor{attributecolor}{rgb}{0.7, 0.1, 0.1} % red
\definecolor{backgroundcolor}{rgb}{1, 1, 1} % light grey
\definecolor{contextcolor}{rgb}{0.9, 0.4, 0.1} %orange
\definecolor{bordercolor}{rgb}{0.5, 0.5, 0.5}
\definecolor{inlinecodecolor}{rgb}{0.92, 0.92, 0.92}

\usepackage{listings}
\def\lstlanguagefiles{lstlean.tex}
\lstset{
  language=lean,
  numbers=none,
  backgroundcolor=\color{gray!10},
  frame=single,
  rulecolor=\color{gray!50},
  frameround=tttt,
  tabsize=2,
  breaklines=true,
  breakatwhitespace=true,
  showstringspaces=false,
  captionpos=b,
  xleftmargin=2em,
  framexleftmargin=1.5em
}